% ---- pg-310 （书页 298） ----

为德宾-沃森(Durbin-Watson)检验.(有关这种检验的详细介绍和临界值表见文献[14].)

\noindent\textbf{例 7.7}\quad (残差图)

例 7.6 中所讨论的计算机数据的三个残差图示于\hyperref[fig:7.3]{图 7.3}. 样本容量 $n=7$ 实在太小, 无法作出明确判断. 不过从图中还是能看出回归假定大体上是成立的. \hfill$\blacksquare$

\begin{mvfigure}
\mvinc{pg-310}{0.153,0.209,0.869,0.311}
\phantomsection\label{fig:7.3}\mvcap{图 7.3\quad 例 7.6 中计算机数据的残差图}
\end{mvfigure}

若对同一些预测变量值能得到若干个响应观测值, 则可以作一个拟合不足的形式检验.(见文献[13]中关于纯误差拟合不足检验的讨论.)

\subsection{杠杆率与影响}

虽然残差分析对评价模型的拟合情况很有用, 但是对回归模型的偏离往往隐藏在拟合过程之中. 例如, 响应变量或解释变量的数据中可能存在``离群值'', 它们对分析会产生很大影响, 但在考察残差图时却不易被发现. 事实上, 这些离群值可能对拟合结果起决定作用.

杠杆率 $h_{jj}$, 即 $\mathbf{H}=\mathbf{Z}(\mathbf{Z}'\mathbf{Z})^{-1}\mathbf{Z}$ 的 $(j,j)$ 对角元, 可以用两种方式解释. 首先, 在解释变量空间中, 杠杆率是第 $j$ 个数据点到其他 $n-1$ 个数据点距离远近的度量. 对于只有一个解释变量 $z$ 的简单线性回归,
\[
h_{jj}=\frac{1}{n}+\frac{(z_j-\bar{z})^2}{\sum\limits_{j=1}^{n}(z_j-\bar{z})^2}
\]
平均杠杆率是 $(r+1)/n$(见习题 7.8).

其次, 杠杆率是单个观测被从拟合中拖开程度的度量. 预测值向量
\[
\hat{\mathbf{y}}=\mathbf{Z}\hat{\boldsymbol{\beta}}=\mathbf{Z}(\mathbf{Z}'\mathbf{Z})^{-1}\mathbf{Z}\mathbf{y}=\mathbf{H}\mathbf{y}
\]
其中第 $j$ 行用观测值将拟合值 $\hat{y}_j$ 表示为
\[
\hat{y}_j=h_{jj}y_j+\sum_{k\neq j}h_{jk}y_k
\]
当所有其他 $y$ 值固定时,
\[
(\hat{y}_j\text{ 的变化})=h_{jj}(y_j\text{ 的变化})
\]
如果杠杆率相对于其他 $h_{jk}$ 比较大, 则 $y_j$ 对预测值 $\hat{y}_j$ 的贡献是主要的.

数据中对分析结果有显著影响的那些观测值称为\textbf{有影响的观测值}. 评估影响大小的方法是根据被考察的观测值被剔除后, 参数估计值向量 $\hat{\boldsymbol{\beta}}$ 的变化程度. 以杠杆率与影响分析统计量为基础的图, 以及它们在诊断性检查中的应用可在文献[3], [5]和[10]中找到.

如果在诊断性检查后没有发现对模型假定的严重背离, 我们就可以对 $\boldsymbol{\beta}$ 和未来的 $Y$ 值作推断, 并在某种程度上能保证不会被误导.
